\chapter{Les archives Bourbaki}\label{ch:bourbaki}

\section{La source}

Les archives de l'Association des collaborateurs de Nicolas Bourbaki (ACNB) sont mises en
ligne par les Archives Henri Poincaré (AHP, Nancy), à
\url{https://archives-bourbaki.ahp-numerique.fr}. Le site, construit sur Omeka, offre une
API publique. Il compte 891 documents : 670 rédactions, 205 circulaires et comptes rendus
de congrès, et quelques autres pièces, soit environ 42\,850 pages. Soixante-dix documents y
sont signalés comme liés au fonds Grothendieck.

\section{Les différences avec Montpellier}

\paragraph{L'accès est simple.} Un PDF par document, à une adresse d'API
(\texttt{api/files?item=<numéro>}), que le serveur laisse afficher dans une autre page :
pas de relais. Le catalogue se lit par l'API, sans gratter de pages HTML.

\paragraph{Ce sont des dactylographies.} La machine les lit bien mieux qu'un brouillon, et
les lots peuvent compter quarante pages. Les PDF n'ont pas de couche de texte, et la
procédure interdit de leur appliquer une reconnaissance de caractères : l'assistant lit
les images, à 110~dpi pour le texte, à 200 ou 300~dpi pour les formules et les ajouts à la
main.

\paragraph{Deux couches, deux mains.} Les symboles mathématiques ont souvent été ajoutés à
la main dans les blancs laissés par la machine à écrire, ou approchés avec les touches
disponibles. La transcription sépare ce que la machine a tapé de ce qu'une main a ajouté,
et n'attribue aucune main à un auteur sans preuve. C'est la question que pose toute
rédaction collective : qui a écrit quoi.

\section{Les droits, et ce qu'on a décidé}

Rien n'est mis en ligne sans l'autorisation de l'AHP et de l'ACNB. Les transcriptions
vivent dans un dépôt privé, et le dépôt public ne garde que ce qui est écrit dans ses propres
mots : le numéro et le titre de chaque rédaction, les énoncés démontrés en Lean, et leur
rédaction. Ce partage entre deux dépôts est l'exemple du chapitre sur les droits de
\emph{Mathématiques outillées}.

L'autorisation a été demandée, et la réponse est venue le 4 novembre 2026 : l'ACNB s'oppose
à l'utilisation de ses archives par des outils d'intelligence artificielle. Le projet a
suspendu tout travail de ce genre sur ces archives, ne lit plus de nouvelle rédaction et ne publie rien de plus.
Quand ce livre est écrit, le site des archives est en accès restreint ; ce qui reste de
public sur Bourbaki dans le dépôt est à décider.

\section{Le bilan}

Dix-sept rédactions de Grothendieck, plus deux longues dactylographies (dont celle de
septembre 1959 qui préfigure les premiers chapitres des \emph{Éléments de géométrie
algébrique}), ont été transcrites localement ; trente-six comptes rendus de congrès ont
été dépouillés. L'article Lean couvre vingt-sept rédactions, de 1950 à 1974. Plusieurs
énoncés dactylographiés s'y révèlent faux, avec un contre-exemple.

\begin{exemple}
Un fonds dactylographié, servi proprement, avec une API, est le cas le plus favorable
techniquement. Le travail porte alors sur ce que la machine lit mal : les ajouts à la main
et les symboles. Et l'accord du détenteur devient la seule vraie question : ici, il a été
refusé, et c'est la leçon qui compte.
\end{exemple}
